\documentclass[10pt,a4paper]{article}
\usepackage[margin=0.5in]{geometry}
\usepackage{amsmath,amssymb}
\usepackage{hyperref}
\usepackage{enumitem}
\usepackage{booktabs}
\usepackage{listings}
\usepackage{xcolor}
\usepackage{titlesec}
\usepackage{parskip}

\titlespacing*{\section}{0pt}{6pt}{2pt}
\titlespacing*{\subsection}{0pt}{4pt}{1pt}
\setlist{itemsep=1pt, topsep=1pt, parsep=0pt}
\setlength{\parskip}{2pt}

\lstset{
    basicstyle=\scriptsize\ttfamily,
    breaklines=true,
    frame=single,
    backgroundcolor=\color{gray!10},
    keywordstyle=\color{blue},
    commentstyle=\color{gray},
    stringstyle=\color{teal},
    xleftmargin=2pt
}

\begin{document}
{\centering \textbf{\large AI Operations (AIOps): Module 1 Assignment Report}\par}
\vspace{4pt}

\section*{Q1: Technical Debt Diagnosis}
\begin{enumerate}[label=(\alph*)]
    \item \textbf{Entanglement (CACE):} changing the delivery-time rounding logic shifted the whole input distribution, which is why the unrelated "favorite restaurants" feature got worse too.
    \item \textbf{Undeclared Consumers:} marketing had been reading the raw output table with no API contract, and nobody on the ML team knew, so a future retrain could silently break their campaign.
    \item \textbf{Config/Glue-Code Debt (Pipeline Jungle):} 14 undocumented shell scripts with no orchestration tool means no single source of truth for run order or configuration, so silent failures are a matter of time.
\end{enumerate}
\textbf{Mitigation for (c):} pin a single \texttt{environment.yml} so the environment is reproducible, swap the loose scripts for an orchestrated \texttt{MLproject} or \texttt{dvc.yaml} where each stage has clearly declared inputs and outputs, and track everything centrally with \texttt{mlflow.log\_param()}, \texttt{log\_metric()}, and \texttt{set\_tag("git\_commit", ...)}.

\section*{Q2: MLflow Experiment Comparison}
I trained an MLP on MNIST across six runs, changing the hidden layer sizes, learning rate, and batch size each time to see what actually moved the needle. \texttt{mlp-wide-256-128} (0.001 lr, batch 128) came out on top at \textbf{97.21\%} val accuracy and 0.9719 F1, probably because the wider two-layer network had more room to pick up stroke and edge patterns than the single-layer baseline. There's no sign of overfitting across the 15 epochs either: train loss dropped from 0.2025 to 0.0069 while val accuracy climbed from 94.02\% to 97.21\% alongside it, with no plateau or reversal. Learning rate mattered the most of everything I tuned. lr=0.0001 underfit badly (94.02\%), lr=0.01 got a bit unstable (95.93\%), and 0.001 landed right in the sweet spot. Architecture width helped a little more on top of that (+0.36\%), but shrinking the batch size to 64 barely moved accuracy (96.89\%) while nearly doubling training time from 24.1s to 40.3s, so that trade-off wasn't worth it.

\begin{center}
\small
\begin{tabular}{lccccc}
\toprule
\textbf{Run} & \textbf{Layers} & \textbf{LR} & \textbf{Batch} & \textbf{Train Loss} & \textbf{Val Acc / F1} \\
\midrule
\texttt{mlp-wide-256-128} & (256,128) & 0.001 & 128 & 0.0069 & \textbf{0.9721} / 0.9719 \\
\texttt{mlp-deep-128-64}  & (128,64)  & 0.001 & 128 & 0.0080 & 0.9705 / 0.9703 \\
\texttt{mlp-small-batch-64} & (128,)  & 0.001 & 64  & 0.0107 & 0.9689 / 0.9687 \\
\texttt{mlp-baseline-128}   & (128,)  & 0.001 & 128 & 0.0229 & 0.9685 / 0.9683 \\
\texttt{mlp-high-lr-0.01}   & (128,)  & 0.010 & 128 & 0.0344 & 0.9593 / 0.9593 \\
\texttt{mlp-low-lr-0.0001}  & (128,)  & 0.0001& 128 & 0.2025 & 0.9402 / 0.9397 \\
\bottomrule
\end{tabular}
\end{center}

\begin{lstlisting}[language=Python]
mlflow.log_param("hidden_layer_sizes", str(hidden_layer_sizes))
mlflow.log_param("learning_rate_init", lr)
mlflow.log_param("batch_size", batch_size)
mlflow.log_param("max_iter", max_iter)
mlflow.log_param("solver", "adam")
mlflow.log_param("random_state", 42)
mlflow.log_metric("train_loss", train_loss)
mlflow.log_metric("val_accuracy", acc)
mlflow.log_metric("val_f1_macro", f1)
mlflow.log_metric("train_time_sec", train_time)
mlflow.set_tag("dataset", "MNIST")
mlflow.sklearn.log_model(model, name="model", serialization_format="cloudpickle")
\end{lstlisting}

\section*{Q3: DVC Data Versioning \& Rollback}
I initialized DVC in the repo and pushed v1 (1800 rows) to an SSH remote\\
(\path{ssh://trader@localhost/home/trader/dvc_storage}), then added the rows from \texttt{new\_labels.zip} to build v2 (2801 rows), re-ran \texttt{dvc add}, and committed that with a clear message. To roll back I ran \texttt{git checkout <v1-commit> -- data.csv.dvc} followed by \texttt{dvc checkout}, which restored v1 exactly, confirmed below:
\begin{lstlisting}
$ git log --oneline --decorate -8     # v1 and v2 are tagged
07cb6f1 (tag: v2) update dataset to v2 2801 lines
8f985a8 (tag: v1) add v1 dataset 1800 rows
$ wc -l data.csv                      # currently on v2
2801 data.csv
$ git checkout v1 -- data.csv.dvc && dvc checkout
M       data.csv
$ wc -l data.csv                      # rolled back to v1
1801 data.csv                         # = 1800 rows + 1 header line
\end{lstlisting}

\section*{Q4: End-to-End Reproducibility Drill}
My partner is Prabhav (DA24B018), and we work in a shared private repo holding only the Q4 materials: \url{https://github.com/Rinkesh-1612/DA3408-Assignment1-Q4}. As Partner A, I trained a RandomForest classifier and logged the params, metrics, seed, a \texttt{git\_commit} tag, and the model artifact to MLflow, then versioned the dataset with DVC and committed the training script together with the \texttt{.dvc} file in a single commit (\texttt{409eef7}). I registered the model as \texttt{q4-digits-classifier} and moved v1 to \textbf{Staging}, giving \texttt{val\_accuracy = 0.9639}. Prabhav reproduces it from the \texttt{partner-a-baseline} tag using only \texttt{git clone}, \texttt{git checkout}, \texttt{dvc checkout}, and \texttt{mamba env create -f environment.yml}, then reruns the script and logs a note recording whether the metric matched within tolerance. The DVC store sits inside that repo, so \texttt{dvc checkout} works on a fresh clone with no credentials. His work goes in under his own commit message, so each half is attributable.

\end{document}
